% =====================================================================
%  Chapitre 4 : Fichiers statiques
% =====================================================================
\chapter{Les fichiers statiques et le design}
\label{ch:statiques}

\chapitreresume{Une page HTML brute est fonctionnelle, mais peu agréable. Dans ce
chapitre, vous ajoutez une feuille de style complète, un logo, un favicon et un
peu de JavaScript. Vous configurez aussi les fichiers média, c'est-à-dire les
images envoyées par les utilisateurs, comme les couvertures de livres.}

\begin{objectifs}
\begin{itemize}
  \item Distinguer fichiers statiques et fichiers média.
  \item Configurer \code{STATIC\_URL}, \code{STATICFILES\_DIRS} et
        \code{STATIC\_ROOT}.
  \item Charger une feuille de style et des images avec \code{\% static \%}.
  \item Écrire un CSS responsive pour le catalogue.
  \item Ajouter un comportement JavaScript simple et accessible.
  \item Configurer l'envoi et l'affichage des couvertures de livres.
  \item Préparer les fichiers statiques pour la production avec
        \code{collectstatic}.
\end{itemize}
\end{objectifs}

\section{Statique ou média ?}

\begin{description}[leftmargin=3.5cm,style=nextline]
  \item[Fichiers statiques] Fichiers livrés avec le site et identiques pour tous
  les visiteurs : CSS, JavaScript, logo, favicon, icônes. Ils sont versionnés avec
  le code.
  \item[Fichiers média] Fichiers envoyés par les utilisateurs : couvertures de
  livres, photos de profil, pièces jointes. Ils changent pendant la vie du site et
  ne sont pas versionnés.
\end{description}

\begin{retenir}
Le CSS, le JavaScript et le logo sont des fichiers statiques. Une couverture de
livre téléversée par un bibliothécaire est un fichier média. Les deux utilisent des
URL, mais pas les mêmes réglages ni le même mode de service en production.
\end{retenir}

\section{Configurer les fichiers statiques}

Ajoutez ces lignes dans \code{settings.py}, après \code{STATIC\_URL} :

\begin{codebox}[python]{gesibiblio/settings.py, extrait}
STATIC_URL = "static/"

# Dossiers où Django cherche les fichiers statiques en développement.
STATICFILES_DIRS = [
    BASE_DIR / "static",
]

# Dossier où collectstatic copie les fichiers pour la production.
STATIC_ROOT = BASE_DIR / "staticfiles"

# Fichiers envoyés par les utilisateurs.
MEDIA_URL = "media/"
MEDIA_ROOT = BASE_DIR / "media"
\end{codebox}

\begin{description}[leftmargin=3.8cm,style=nextline]
  \item[\code{STATIC\_URL}] Préfixe des URL statiques. Avec \code{"static/"},
  le fichier \code{static/css/style.css} est servi à l'adresse
  \code{/static/css/style.css}.
  \item[\code{STATICFILES\_DIRS}] Dossiers supplémentaires scannés par Django.
  Ici, le dossier \code{static/} à la racine du projet.
  \item[\code{STATIC\_ROOT}] Dossier de destination de \code{collectstatic}. Il ne
  faut pas le versionner.
  \item[\code{MEDIA\_URL}] Préfixe des URL des fichiers envoyés.
  \item[\code{MEDIA\_ROOT}] Dossier de stockage des fichiers envoyés.
\end{description}

\begin{attention}
Ne confondez pas \code{STATICFILES\_DIRS} et \code{STATIC\_ROOT}.
\code{STATICFILES\_DIRS} indique où vous écrivez vos fichiers.
\code{STATIC\_ROOT} indique où \code{collectstatic} les rassemble pour la
production. Écrire directement dans \code{STATIC\_ROOT} est une erreur.
\end{attention}

Créez l'arborescence suivante :

\begin{codebox}[tree]{Arborescence des fichiers statiques}
static/
|-- css/
|   `-- style.css
|-- js/
|   `-- app.js
`-- img/
    |-- logo-gesibiblio.svg
    `-- favicon.svg

media/
`-- couvertures/
\end{codebox}

\section{Charger un fichier statique}

Dans un gabarit, commencez par charger la bibliothèque \code{static}, puis utilisez
la balise \code{static} :

\begin{codebox}[html]{templates/base.html}
{% load static %}
<!DOCTYPE html>
<html lang="fr">
<head>
  <meta charset="utf-8">
  <meta name="viewport" content="width=device-width, initial-scale=1">
  <title>{% block titre %}GesiBiblio{% endblock %}</title>
  <link rel="stylesheet" href="{% static 'css/style.css' %}">
  <link rel="icon" href="{% static 'img/favicon.svg' %}" type="image/svg+xml">
</head>
<body>
  {% include "partials/navbar.html" %}
  {% include "partials/messages.html" %}

  <main class="conteneur">
    {% block contenu %}{% endblock %}
  </main>

  {% include "partials/footer.html" %}
  <script src="{% static 'js/app.js' %}" defer></script>
</body>
</html>
\end{codebox}

\begin{attention}
\code{\{\% load static \%\}} doit apparaître dans chaque gabarit qui utilise
\code{\{\% static \%\}}. Le charger dans \code{base.html} ne le rend pas disponible
dans les gabarits enfants : chaque fichier a son propre contexte de bibliothèques.
\end{attention}

\section{Le logo et le favicon}

Le logo de l'application est un fichier SVG. Il est net à toutes les tailles et
pèse quelques centaines d'octets. Créez \code{static/img/logo-gesibiblio.svg} :

\begin{codebox}[html]{static/img/logo-gesibiblio.svg}
<svg xmlns="http://www.w3.org/2000/svg" viewBox="0 0 64 64"
     role="img" aria-label="GesiBiblio">
  <rect width="64" height="64" rx="16" fill="#1F3A5F"/>
  <path d="M15 17h13c3 0 5 2 5 5v26c0-3-2-5-5-5H15z" fill="#E9C46A"/>
  <path d="M49 17H36c-3 0-5 2-5 5v26c0-3 2-5 5-5h13z" fill="#2A9D8F"/>
  <path d="M32 22v26" stroke="#FFFFFF" stroke-width="2"/>
</svg>
\end{codebox}

Le favicon est une version simplifiée, affichée dans l'onglet du navigateur :

\begin{codebox}[html]{static/img/favicon.svg}
<svg xmlns="http://www.w3.org/2000/svg" viewBox="0 0 64 64">
  <rect width="64" height="64" rx="14" fill="#1F3A5F"/>
  <path d="M16 18h12c3 0 5 2 5 5v24c0-3-2-5-5-5H16z" fill="#E9C46A"/>
  <path d="M48 18H36c-3 0-5 2-5 5v24c0-3 2-5 5-5h12z" fill="#2A9D8F"/>
</svg>
\end{codebox}

\section{La feuille de style}

Nous allons construire une interface sobre et lisible, inspirée des sites de
bibliothèque modernes. Les couleurs reprennent la charte du document : bleu nuit,
vert d'eau, ambre et corail.

\subsection{Variables et base}

\begin{codebox}[css]{static/css/style.css, partie 1}
/* Variables de la charte graphique */
:root {
  --navy: #1f3a5f;
  --navy-2: #2e5077;
  --teal: #2a9d8f;
  --teal-2: #7fc8be;
  --amber: #e9c46a;
  --coral: #e76f51;
  --grey: #5c677d;
  --bg: #f6f8fa;
  --white: #ffffff;
  --border: #dfe6ee;
  --radius: 14px;
  --shadow: 0 12px 30px rgba(31, 58, 95, 0.08);
}

*, *::before, *::after {
  box-sizing: border-box;
}

html {
  scroll-behavior: smooth;
}

body {
  margin: 0;
  font-family: system-ui, -apple-system, "Segoe UI", Roboto,
               "Helvetica Neue", Arial, sans-serif;
  color: var(--navy);
  background: var(--bg);
  line-height: 1.65;
}

img {
  max-width: 100%;
  display: block;
}

a {
  color: var(--teal);
  text-decoration: none;
}

a:hover {
  text-decoration: underline;
}

h1, h2, h3 {
  line-height: 1.2;
  color: var(--navy);
  margin-top: 0;
}

.conteneur {
  width: min(1120px, 100% - 2.5rem);
  margin-inline: auto;
}
\end{codebox}

\subsection{En-tête, navigation et boutons}

\begin{codebox}[css]{static/css/style.css, partie 2}
.entete {
  position: sticky;
  top: 0;
  z-index: 20;
  background: rgba(255, 255, 255, 0.96);
  border-bottom: 1px solid var(--border);
  backdrop-filter: blur(8px);
}

.barre {
  display: flex;
  align-items: center;
  justify-content: space-between;
  min-height: 72px;
  gap: 1rem;
  position: relative;
}

.marque {
  display: inline-flex;
  align-items: center;
  gap: 0.65rem;
  font-weight: 800;
  font-size: 1.2rem;
  color: var(--navy);
}

.marque:hover {
  text-decoration: none;
}

.menu {
  display: flex;
  align-items: center;
  gap: 1.4rem;
  list-style: none;
  margin: 0;
  padding: 0;
}

.menu a {
  color: var(--navy);
  font-weight: 600;
  font-size: 0.95rem;
}

.menu a:hover {
  color: var(--teal);
}

.burger {
  display: none;
  border: 1px solid var(--border);
  background: var(--white);
  color: var(--navy);
  border-radius: 10px;
  padding: 0.5rem 0.8rem;
  font: inherit;
  font-weight: 600;
  cursor: pointer;
}

.bouton {
  display: inline-flex;
  align-items: center;
  justify-content: center;
  gap: 0.4rem;
  padding: 0.7rem 1.15rem;
  border-radius: 10px;
  border: 1px solid var(--border);
  background: var(--white);
  color: var(--navy);
  font-weight: 700;
  font-size: 0.95rem;
  cursor: pointer;
  transition: transform 0.15s ease, box-shadow 0.15s ease;
}

.bouton:hover {
  text-decoration: none;
  transform: translateY(-1px);
  box-shadow: var(--shadow);
}

.bouton-principal {
  background: var(--teal);
  border-color: var(--teal);
  color: var(--white);
}

.bouton-danger {
  background: var(--coral);
  border-color: var(--coral);
  color: var(--white);
}

.lien {
  background: none;
  border: 0;
  color: var(--teal);
  font: inherit;
  cursor: pointer;
  padding: 0;
}
\end{codebox}

\subsection{Héros, statistiques et grille de livres}

\begin{codebox}[css]{static/css/style.css, partie 3}
.hero {
  display: grid;
  grid-template-columns: 1.4fr 1fr;
  align-items: center;
  gap: 2rem;
  margin: 2.5rem 0;
  padding: 2.5rem;
  border-radius: 24px;
  background: linear-gradient(135deg, #eaf6f4 0%, #fdf6e3 100%);
  border: 1px solid var(--border);
}

.hero h1 {
  font-size: clamp(2rem, 4vw, 3.2rem);
  margin-bottom: 0.6rem;
}

.hero p {
  font-size: 1.1rem;
  color: var(--grey);
  max-width: 52ch;
}

.surtitre {
  text-transform: uppercase;
  letter-spacing: 0.12em;
  font-size: 0.78rem;
  font-weight: 800;
  color: var(--teal);
}

.statistiques {
  display: grid;
  grid-template-columns: repeat(3, 1fr);
  gap: 1rem;
  margin: 2rem 0;
}

.statistiques div {
  background: var(--white);
  border: 1px solid var(--border);
  border-radius: var(--radius);
  padding: 1.2rem;
  text-align: center;
}

.statistiques strong {
  display: block;
  font-size: 2rem;
  color: var(--teal);
}

.grille-livres {
  display: grid;
  grid-template-columns: repeat(auto-fill, minmax(220px, 1fr));
  gap: 1.4rem;
  margin: 1.5rem 0 3rem;
}

.carte-livre {
  background: var(--white);
  border: 1px solid var(--border);
  border-radius: var(--radius);
  overflow: hidden;
  display: flex;
  flex-direction: column;
  transition: transform 0.2s ease, box-shadow 0.2s ease;
}

.carte-livre:hover {
  transform: translateY(-4px);
  box-shadow: var(--shadow);
}

.carte-livre img,
.couverture-vide {
  width: 100%;
  aspect-ratio: 2 / 3;
  object-fit: cover;
}

.couverture-vide {
  display: grid;
  place-items: center;
  background: linear-gradient(160deg, var(--navy), var(--teal));
  color: var(--white);
  font-size: 3rem;
  font-weight: 800;
}

.carte-livre h3 {
  font-size: 1.05rem;
  margin: 0.9rem 0.9rem 0.2rem;
}

.carte-livre .auteur {
  margin: 0 0.9rem;
  color: var(--grey);
  font-size: 0.9rem;
}
\end{codebox}

\subsection{Badges, fiches, formulaires et tableaux}

\begin{codebox}[css]{static/css/style.css, partie 4}
.badge {
  display: inline-block;
  align-self: flex-start;
  margin: 0.7rem 0.9rem 0.9rem;
  padding: 0.2rem 0.6rem;
  border-radius: 999px;
  font-size: 0.75rem;
  font-weight: 800;
}

.badge-disponible {
  background: #dff3ee;
  color: #1d6b5f;
}

.badge-indisponible {
  background: #fde7e2;
  color: #a33a22;
}

.badge-retard {
  background: #fde7e2;
  color: #a33a22;
}

.fiche-livre {
  display: grid;
  grid-template-columns: 300px 1fr;
  gap: 2.2rem;
  margin: 2.5rem 0;
  align-items: start;
}

.fiche-livre img {
  border-radius: var(--radius);
  box-shadow: var(--shadow);
}

.fiche-livre dl {
  display: grid;
  grid-template-columns: 9rem 1fr;
  gap: 0.4rem 1rem;
  margin: 1.2rem 0;
}

.fiche-livre dt {
  font-weight: 800;
  color: var(--navy);
}

.fiche-livre dd {
  margin: 0;
  color: var(--grey);
}

.formulaire {
  background: var(--white);
  border: 1px solid var(--border);
  border-radius: var(--radius);
  padding: 1.6rem;
  max-width: 720px;
  margin: 1.5rem 0 3rem;
}

.champ {
  margin-bottom: 1.1rem;
}

.champ label {
  display: block;
  font-weight: 700;
  margin-bottom: 0.35rem;
}

.champ input,
.champ select,
.champ textarea {
  width: 100%;
  padding: 0.65rem 0.75rem;
  border: 1px solid var(--border);
  border-radius: 9px;
  font: inherit;
  color: var(--navy);
  background: var(--white);
}

.champ input:focus,
.champ select:focus,
.champ textarea:focus {
  outline: 2px solid var(--teal-2);
  border-color: var(--teal);
}

.aide {
  margin: 0.3rem 0 0;
  font-size: 0.82rem;
  color: var(--grey);
}

.erreur {
  margin: 0.3rem 0 0;
  font-size: 0.85rem;
  color: var(--coral);
  font-weight: 600;
}

.tableau {
  width: 100%;
  border-collapse: collapse;
  background: var(--white);
  border: 1px solid var(--border);
  border-radius: var(--radius);
  overflow: hidden;
  margin: 1.5rem 0 3rem;
}

.tableau th,
.tableau td {
  padding: 0.8rem 1rem;
  text-align: left;
  border-bottom: 1px solid var(--border);
  font-size: 0.92rem;
}

.tableau th {
  background: #eef3f8;
  color: var(--navy);
  font-size: 0.82rem;
  text-transform: uppercase;
  letter-spacing: 0.04em;
}

.tableau tr:last-child td {
  border-bottom: 0;
}

.tableau tbody tr:nth-child(even) {
  background: #fbfcfe;
}
\end{codebox}

\subsection{Messages, pagination, pied de page et responsive}

\begin{codebox}[css]{static/css/style.css, partie 5}
.messages {
  margin-top: 1rem;
}

.message {
  padding: 0.8rem 1rem;
  border-radius: 10px;
  border: 1px solid var(--border);
  background: var(--white);
  font-weight: 600;
  transition: opacity 0.4s ease, transform 0.4s ease;
}

.message.success {
  border-color: var(--teal);
  background: #e8f7f3;
  color: #1d6b5f;
}

.message.error {
  border-color: var(--coral);
  background: #fdece8;
  color: #a33a22;
}

.message.warning {
  border-color: var(--amber);
  background: #fdf6e3;
  color: #7a5b12;
}

.message.disparition {
  opacity: 0;
  transform: translateY(-6px);
}

.pagination {
  display: flex;
  align-items: center;
  justify-content: center;
  gap: 1rem;
  margin: 2rem 0 3rem;
  font-weight: 700;
}

.pied {
  margin-top: 4rem;
  padding: 2rem 0;
  background: var(--navy);
  color: #dce6f0;
  font-size: 0.9rem;
}

.pied p {
  margin: 0.2rem 0;
}

@media (max-width: 820px) {
  .burger {
    display: inline-flex;
  }

  .menu {
    display: none;
    position: absolute;
    top: 100%;
    left: 0;
    right: 0;
    flex-direction: column;
    align-items: flex-start;
    padding: 1rem 1.25rem;
    background: var(--white);
    border-bottom: 1px solid var(--border);
    box-shadow: var(--shadow);
  }

  .menu.ouvert {
    display: flex;
  }

  .hero,
  .fiche-livre {
    grid-template-columns: 1fr;
    padding: 1.5rem;
  }

  .statistiques {
    grid-template-columns: 1fr;
  }

  .fiche-livre dl {
    grid-template-columns: 1fr;
  }
}
\end{codebox}

\section{Le JavaScript}

Le JavaScript reste minimal : ouvrir et fermer le menu sur mobile, et faire
disparaître les messages après quelques secondes.

\begin{codebox}[javascript]{static/js/app.js}
const burger = document.querySelector(".burger");
const menu = document.querySelector("#menu-principal");

if (burger && menu) {
  burger.addEventListener("click", () => {
    const ouvert = menu.classList.toggle("ouvert");
    burger.setAttribute("aria-expanded", String(ouvert));
  });
}

document.querySelectorAll(".message").forEach((message) => {
  window.setTimeout(() => {
    message.classList.add("disparition");
  }, 5000);
});
\end{codebox}

\begin{bonne}
Un site doit rester utilisable sans JavaScript. Ici, la navigation fonctionne même
si le script ne se charge pas : le menu est visible par défaut sur ordinateur et
les liens restent accessibles. Le JavaScript améliore l'expérience, il ne porte
pas toute la logique.
\end{bonne}

\section{Les fichiers média}

Les couvertures de livres sont envoyées par les utilisateurs. Elles ne sont pas
statiques : elles vivent dans \code{MEDIA\_ROOT} et sont servies via
\code{MEDIA\_URL}.

\subsection{Configuration}

\begin{codebox}[python]{gesibiblio/settings.py, extrait}
MEDIA_URL = "media/"
MEDIA_ROOT = BASE_DIR / "media"
\end{codebox}

En développement, Django doit servir ces fichiers. Ajoutez les lignes suivantes à
la fin du fichier \code{urls.py} du projet :

\begin{codebox}[python]{gesibiblio/urls.py}
from django.conf import settings
from django.conf.urls.static import static
from django.contrib import admin
from django.urls import include, path

urlpatterns = [
    path("admin/", admin.site.urls),
    path("", include("catalogue.urls")),
]

if settings.DEBUG:
    urlpatterns += static(
        settings.MEDIA_URL,
        document_root=settings.MEDIA_ROOT,
    )
\end{codebox}

\begin{attention}
\code{static()} ne fonctionne qu'en développement, lorsque \code{DEBUG} vaut
\code{True}. En production, Nginx ou un service de stockage doit servir les
fichiers média. Ne laissez jamais Django servir des fichiers envoyés par des
utilisateurs en production.
\end{attention}

\subsection{Afficher une couverture}

Dans un gabarit, testez d'abord la présence de l'image :

\begin{codebox}[html]{templates/catalogue/livres/carte.html}
{% if livre.couverture %}
  <img src="{{ livre.couverture.url }}"
       alt="Couverture de {{ livre.titre }}">
{% else %}
  <div class="couverture-vide" aria-hidden="true">
    {{ livre.titre|first }}
  </div>
{% endif %}
\end{codebox}

\begin{note}
\code{livre.couverture} est un objet \code{ImageField}. \code{.url} renvoie
l'adresse publique du fichier, par exemple \code{/media/couvertures/le-petit-prince.jpg}.
L'image n'est accessible que si \code{MEDIA\_URL} et \code{MEDIA\_ROOT} sont
correctement configurés.
\end{note}

\section{Préparer la production}

En développement, Django sert les fichiers statiques automatiquement. En
production, ce n'est pas le cas. La commande \code{collectstatic} rassemble tous
les fichiers dans \code{STATIC\_ROOT} :

\begin{terminal}
python manage.py collectstatic
\end{terminal}

Le serveur web (Nginx, Apache) ou un service comme WhiteNoise sert ensuite ce
dossier. Une configuration typique avec WhiteNoise :

\begin{terminal}
python -m pip install whitenoise
\end{terminal}

\begin{codebox}[python]{gesibiblio/settings.py, extrait}
MIDDLEWARE = [
    "django.middleware.security.SecurityMiddleware",
    "whitenoise.middleware.WhiteNoiseMiddleware",
    # ... autres middlewares
]

STORAGES = {
    "default": {
        "BACKEND": "django.core.files.storage.FileSystemStorage",
    },
    "staticfiles": {
        "BACKEND": "whitenoise.storage.CompressedManifestStaticFilesStorage",
    },
}
\end{codebox}

\begin{bonne}
En production, utilisez un nom de fichier avec un hash, comme
\code{style.3f2a1c.css}, pour que les navigateurs mettent à jour le cache. Le
\code{CompressedManifestStaticFilesStorage} de WhiteNoise s'en charge. En
développement, gardez la configuration simple.
\end{bonne}

\section{Exercices}

\begin{exercice}[ : changer la charte]
Modifiez les variables CSS pour adopter une palette différente : un bleu plus
clair, un vert plus doux, un accent orange. Rechargez la page et observez les
zones concernées.
\end{exercice}

\begin{correction}
Modifiez uniquement le bloc \code{:root}. Par exemple :

\begin{codebox}[css]{static/css/style.css, extrait}
:root {
  --navy: #17324d;
  --teal: #3aa6a0;
  --amber: #f2c14e;
  --coral: #ef6f6c;
  --bg: #f7f9fb;
}
\end{codebox}

Toutes les règles qui utilisent \code{var(--teal)} changent automatiquement. C'est
tout l'intérêt des variables CSS : une seule source de vérité.
\end{correction}

\begin{exercice}[ : page 404 personnalisée]
Créez un gabarit \code{templates/404.html} qui étend \code{base.html} et affiche un
message sympathique avec un lien vers l'accueil. Testez-le en visitant une URL
inexistante avec \code{DEBUG = False}.
\end{exercice}

\begin{correction}
\begin{codebox}[html]{templates/404.html}
{% extends "base.html" %}

{% block titre %}Page introuvable | GesiBiblio{% endblock %}

{% block contenu %}
  <section class="hero">
    <div>
      <p class="surtitre">Erreur 404</p>
      <h1>Cette page est introuvable.</h1>
      <p>Le livre que vous cherchez a peut-être changé de rayon.</p>
      <a class="bouton bouton-principal"
         href="{% url 'catalogue:accueil' %}">Retour à l'accueil</a>
    </div>
  </section>
{% endblock %}
\end{codebox}

Pour tester, passez \code{DEBUG = False} et ajoutez \code{localhost} et
\code{127.0.0.1} à \code{ALLOWED\_HOSTS}, puis redémarrez le serveur. Django
affiche alors \code{404.html} au lieu de la page de débogage.
\end{correction}

\begin{exercice}[ : favicon]
Ajoutez le favicon dans \code{base.html} et vérifiez qu'il apparaît dans l'onglet
du navigateur. Si l'ancien favicon reste affiché, videz le cache ou ouvrez une
fenêtre de navigation privée.
\end{exercice}

\begin{correction}
La ligne est déjà présente dans le \code{base.html} complet :

\begin{codebox}[html]{templates/base.html, extrait}
<link rel="icon" href="{% static 'img/favicon.svg' %}" type="image/svg+xml">
\end{codebox}

Vérifiez que le fichier \code{static/img/favicon.svg} existe et que
\code{STATICFILES\_DIRS} pointe bien vers \code{BASE\_DIR / "static"}.
\end{correction}

\begin{retenir}
\begin{itemize}
  \item Les fichiers statiques sont livrés avec le site ; les médias sont envoyés
        par les utilisateurs.
  \item \code{STATICFILES\_DIRS} indique où écrire, \code{STATIC\_ROOT} où
        collecter pour la production.
  \item \code{\{\% static \%\}} construit les URL des fichiers statiques.
  \item \code{MEDIA\_URL} et \code{MEDIA\_ROOT} servent les fichiers envoyés.
  \item \code{collectstatic} prépare les fichiers pour la production.
  \item Un design responsive repose sur des variables, une grille souple et des
        points de rupture simples.
\end{itemize}
\end{retenir}
